光伏出力与负载需求在时间上的不匹配，使微网购电同时受到供需变化、电价波动和储能状态的影响。本文以购电费用最小为目标，建立单日联合优化模型，并在波动电价条件下构建预测、调度与实际结算相衔接的滚动模型。

\textbf{针对问题一}，将全天划分为144个十分钟区间，以145个边界储电量连接各时段决策。按照母线侧电量建立供需平衡，取充、放电效率各为0.9，通过二进制变量限制同时充放电，并要求日末储电量恢复至日初水平，得到混合整数线性规划模型。采用HiGHS求解，全天购电量为\textbf{59\,482.699 kWh}，购电费用为\textbf{35\,126.949元}，55\,482.836 kWh光伏预测电量全部消纳，初末储电量均为6\,000 kWh。求解器报告相对最优间隙为0；独立复算的最大供需平衡残差和储能递推残差分别为$1.14\times10^{-13}$和$9.09\times10^{-13}$ kWh。

\textbf{针对问题二}，以可见历史构造负载周周期--残差修正预测和光伏日滞后预测，在每日0:00冻结正常购电合同，并以24小时滚动模型安排十分钟储能动作。实际回放覆盖2025年2月至12月的48\,096个区间，正常合同费、紧急购电费和总费用分别为\textbf{1350.0196万元}、\textbf{168.5776万元}和\textbf{1518.5972万元}。实际总费用较1500万元目标高18.5972万元，紧急购电费较100万元目标高68.5776万元。紧急购电遍及333个回放日，59.35\%集中于10:00--16:00；配对预测表明，午间光伏高估与日前合同冻结共同将净负荷误差转化为高价补缺。本文据实际账单报告该未达标结果，不以预测目标代替结算结果。

\textbf{针对问题四}，区分每日合同冻结与允许按预报调整合同两类条件，以历史基线和残差修正生成负载、电价预测，按历史误差组合已发布光伏预报。利用过去28日已实现的联合预测误差构造低、名义、高净需求情景，在各情景共享合同和充放电动作的条件下滚动优化，并按实际供需反馈和交易记录结算。情景方案与点预测基准均完成2025年2月至12月的334日连续回放，实际末态均为6\,000 kWh。两类情景方案费用分别为\textbf{1591.9083万元}和\textbf{1540.2971万元}，相对同条件点预测基准减少\textbf{7.3466\%}和\textbf{2.9194\%}；紧急购电量分别减少45.7195、28.5367万kWh。

计算结果表明，预测误差在合同冻结条件下会转化为高价紧急购电；情景采购则通过增加正常合同购电、减少高价缺口补电降低总费，但合同余电与弃光增加。上述节费不来自点预测精度提高，且仍需独立年份验证。模型未计电池寿命费用，问题三的独立结果尚待补充。
